\chapter{Conclusion \& Future Scope}

\section{Project Summary}
Recall demonstrates that a small, dependency-light application can turn arbitrary topics and study documents into complete, gradable quizzes. Source text is retrieved from the web or extracted from files, wrapped in a schema-constrained prompt, sent to a cloud or local language model, parsed, validated and delivered through an exam-style interface with AI grading and a code runner. On top of this, an adaptive learning layer (Chapter~\ref{chap:adaptive}) tracks mastery per concept, schedules reviews, checks each question against its source and plans revision for an exam date, almost entirely in the browser. The key engineering ideas are: grounding questions in visible sources, never trusting raw model output, always providing a fallback path, and keeping the server stateless.

\section{Current System Limitations}
\begin{itemize}
\item Only Python code can be executed, and the runner is not isolated from the host.
\item No user accounts; history is local to one browser.
\item No optical character recognition for scanned documents.
\item Question quality depends on the chosen model and, for topic quizzes, on the quality of retrieved sources.
\item The mastery model, the hint and response-time rules and the retention curve use fixed default parameters that have not been fitted to data from real learners.
\item The evidence check is lexical and cannot judge whether a quote really justifies an answer.
\item On the free cloud tier a request with the largest question counts can exceed the per-minute token limit.
\end{itemize}

\section{Objectives Traceability Matrix}
\begin{table}[H]
\centering
\begin{tabularx}{\textwidth}{@{} p{1.3cm} X p{3.6cm} @{}}
\toprule
\textbf{Obj.} & \textbf{Objective (Section 1.3)} & \textbf{Where realised} \\
\midrule
1 & Topic quizzes grounded in web sources & Ch.~7, \texttt{research.py} \\
2 & Quizzes from uploaded documents and pasted text & Ch.~7, \texttt{file\_loader.py} \\
3 & Three question types and three difficulties & Ch.~5, Ch.~6 \\
4 & Strict JSON schema, validation and retry & Ch.~6, \texttt{quiz\_engine.py} \\
5 & Cloud engine with offline local fallback & Ch.~6 \\
6 & AI grading and code execution & Ch.~6, Ch.~8 \\
7 & Practice and Exam modes, history, export & Ch.~9 \\
\bottomrule
\end{tabularx}
\caption{Objectives traceability matrix}
\label{tab:trace}
\end{table}

\section{Future Architectural Enhancements}
\begin{itemize}
\item \textbf{Containerised sandbox} supporting Java, C, C++ and JavaScript with resource limits.
\item \textbf{Continuous integration} that runs the existing \texttt{pytest} and Node test suites on every change.
\item \textbf{Accounts and cloud sync} for history and learning data, so that mastery and review schedules follow the learner across devices.
\item \textbf{Fitting the learning model} (BKT parameters, retention curve) to answer data from real learners.
\item \textbf{OCR} for scanned PDFs and images of handwritten notes.
\item \textbf{Semantic question checks}, such as a second model pass that verifies each answer against its source quote, to complement the lexical evidence check.
\item \textbf{Teacher mode} for sharing a generated quiz through a link and collecting results.
\item \textbf{Streaming generation} so that questions appear as soon as they are produced.
\end{itemize}
